% GENERATED FILE. Edit canonical lessons, then run npm run generate:books.
% canonical-source-hash: fnv1a64:5130c8465fd5cba3
% canonical-lessons: MR-C343-mudda, MR-C343-tarka, MR-C343-yuktivad, MR-C343-mandne, MR-C343-purava

\chapter{Making a Case}
\label{ch:mr-making-a-case}
\hlchaptermodality{343}

\begin{chapteropening}
\textbf{By the end of this chapter:} \emph{I can put forward a view and back it with reasons.}

Complete the last lesson of chapter 343: I can put forward a view and back it with reasons.
\end{chapteropening}

\section[\texorpdfstring{\mr{मुद्दा} (muddā)}{मुद्दा (muddā)}]{\mr{मुद्दा} (muddā) — a point at issue; a point in an argument}

\label{lesson:MR-C343-mudda}



\begin{quote}
Before the new one: say the Marathi for an explanation, then the Marathi for to come about.
\end{quote}

\subsection*{You'll want to know: \mr{मुद्दा}}
\textbf{\mr{मुद्दा}} — \emph{muddā} — \textquotedblleft{}a point at issue; a point in an argument\textquotedblright{}.

\begin{quote}
\mr{हा} \mr{मुद्दा} \mr{फार} \mr{महत्त्वाचा} \mr{आहे}.
\end{quote}

\begin{quote}
\emph{hā muddā phār mahattvācā āhe.}
\end{quote}

\begin{quote}
This point is very important.
\end{quote}

Masculine. Plural: \mr{मुद्दे}, \emph{mudde}.

\mr{मुद्द्यावर} \mr{येणे}, to come to the point.

\begin{cousinweb}
\textbf{\mr{मुद्दा}} (\emph{muddā}) is a loan from Persian, usually traced to Arabic \emph{muddaʿā}, \textquotedblleft{}a thing claimed or alleged\textquotedblright{}, from the same Arabic root as \emph{daʿwā}, \textquotedblleft{}a claim\textquotedblright{}; in Marathi it is the point under discussion.
\end{cousinweb}

\begin{grammarlens}[title={What you've built}]
One more word for giving reasons and opinions.
\end{grammarlens}

\subsection*{Your turn}
\begin{itemize}
\raggedright
  \item \emph{Say it:} \emph{muddā}
  \item \emph{Say it:} \emph{muddā}, once more
  \item \emph{Say it:} \emph{hā muddā phār mahattvācā āhe.} — This point is very important.
  \item \emph{Recall:} say the Marathi for an explanation, then the Marathi for to come about, then say \emph{muddā} again
\end{itemize}

\begin{tcolorbox}[breakable,colback=teal!4,colframe=teal!35!black,title={Before you move on}]
What does \mr{मुद्दा} mean? (\textquotedblleft{}A point at issue; a point in an argument\textquotedblright{}.) Say it once more.
\end{tcolorbox}
\section[\texorpdfstring{\mr{तर्क} (tarka)}{तर्क (tarka)}]{\mr{तर्क} (tarka) — logic; reasoning}

\label{lesson:MR-C343-tarka}



\begin{quote}
Before the new one: say the Marathi for to come about, then the Marathi for a point at issue.
\end{quote}

\subsection*{You'll want to know: \mr{तर्क}}
\textbf{\mr{तर्क}} — \emph{tarka} — \textquotedblleft{}logic; reasoning\textquotedblright{}.

\begin{quote}
\mr{त्याच्या} \mr{बोलण्यात} \mr{काहीच} \mr{तर्क} \mr{नव्हता}.
\end{quote}

\begin{quote}
\emph{tyācyā bolaṇyāt kāhīc tarka navhatā.}
\end{quote}

\begin{quote}
There was no logic at all in what he said.
\end{quote}

Masculine.

\mr{तर्क} \mr{लावणे}, to reason something out, to guess.

\begin{cousinweb}
\textbf{\mr{तर्क}} (\emph{tarka}) is a tatsama, taken from Sanskrit \emph{tarka}, \textquotedblleft{}conjecture, reasoning\textquotedblright{} (Monier-Williams). CDIAL lists the same head (CDIAL 5714).
\end{cousinweb}

\begin{grammarlens}[title={What you've built}]
One more word for giving reasons and opinions.
\end{grammarlens}

\subsection*{Your turn}
\begin{itemize}
\raggedright
  \item \emph{Say it:} \emph{tarka}
  \item \emph{Say it:} \emph{tarka}, once more
  \item \emph{Say it:} \emph{tyācyā bolaṇyāt kāhīc tarka navhatā.} — There was no logic at all in what he said.
  \item \emph{Recall:} say the Marathi for to come about, then the Marathi for a point at issue, then say \emph{tarka} again
\end{itemize}

\begin{tcolorbox}[breakable,colback=teal!4,colframe=teal!35!black,title={Before you move on}]
What does \mr{तर्क} mean? (\textquotedblleft{}Logic; reasoning\textquotedblright{}.) Say it once more.
\end{tcolorbox}
\section[\texorpdfstring{\mr{युक्तिवाद} (yuktivād)}{युक्तिवाद (yuktivād)}]{\mr{युक्तिवाद} (yuktivād) — a reasoned argument}

\label{lesson:MR-C343-yuktivad}



\begin{quote}
Before the new one: say the Marathi for a point at issue, then the Marathi for logic.
\end{quote}

\subsection*{You'll want to know: \mr{युक्तिवाद}}
\textbf{\mr{युक्तिवाद}} — \emph{yuktivād} — \textquotedblleft{}a reasoned argument\textquotedblright{}.

\begin{quote}
\mr{त्याचा} \mr{युक्तिवाद} \mr{ऐकून} \mr{सगळे} \mr{गप्प} \mr{बसले}.
\end{quote}

\begin{quote}
\emph{tyācā yuktivād aikūn sagḷe gappa basle.}
\end{quote}

\begin{quote}
Hearing his argument, everyone fell silent.
\end{quote}

Masculine.

\mr{युक्तिवाद} \mr{करणे}, to argue a case; a quarrel is \mr{भांडण}.

\begin{cousinweb}
\textbf{\mr{युक्तिवाद}} (\emph{yuktivād}) is built from two tatsamas, Sanskrit \emph{yukti}, \textquotedblleft{}contrivance, reasoning\textquotedblright{}, and \emph{vāda}, \textquotedblleft{}speech, thesis\textquotedblright{} (Monier-Williams).
\end{cousinweb}

\begin{grammarlens}[title={What you've built}]
One more word for giving reasons and opinions.
\end{grammarlens}

\subsection*{Your turn}
\begin{itemize}
\raggedright
  \item \emph{Say it:} \emph{yuktivād}
  \item \emph{Say it:} \emph{yuktivād}, once more
  \item \emph{Say it:} \emph{tyācā yuktivād aikūn sagḷe gappa basle.} — Hearing his argument, everyone fell silent.
  \item \emph{Recall:} say the Marathi for a point at issue, then the Marathi for logic, then say \emph{yuktivād} again
\end{itemize}

\begin{tcolorbox}[breakable,colback=teal!4,colframe=teal!35!black,title={Before you move on}]
What does \mr{युक्तिवाद} mean? (\textquotedblleft{}A reasoned argument\textquotedblright{}.) Say it once more.
\end{tcolorbox}
\section[\texorpdfstring{\mr{मांडणे} (māṇḍṇe)}{मांडणे (māṇḍṇe)}]{\mr{मांडणे} (māṇḍṇe) — to put forward; to set out}

\label{lesson:MR-C343-mandne}



\begin{quote}
Before the new one: say the Marathi for logic, then the Marathi for a reasoned argument.
\end{quote}

\subsection*{You'll want to know: \mr{मांडणे}}
\textbf{\mr{मांडणे}} — \emph{māṇḍṇe} — \textquotedblleft{}to put forward; to set out\textquotedblright{}.

\begin{quote}
\mr{मीराने} \mr{बैठकीत} \mr{आपलं} \mr{मत} \mr{मांडलं}.
\end{quote}

\begin{quote}
\emph{mīrāne baiṭhkīt āplaṁ mat māṇḍlaṁ.}
\end{quote}

\begin{quote}
Meera put forward her view at the meeting.
\end{quote}

Past, with -\mr{ने} (he …): \mr{त्याने} \mr{मुद्दा} \mr{मांडला}, \emph{tyāne muddā māṇḍlā}.

\mr{मत} \mr{मांडणे}, to state one's view; also to set things out in order, as on a table.

\begin{cousinweb}
\textbf{\mr{मांडणे}} (\emph{māṇḍṇe}) is a Marathi verb whose origin is not settled; its first sense is to set things out in order, and from there to set out a view.
\end{cousinweb}

\begin{grammarlens}[title={What you've built}]
One more word for giving reasons and opinions.
\end{grammarlens}

\subsection*{Your turn}
\begin{itemize}
\raggedright
  \item \emph{Say it:} \emph{māṇḍṇe}
  \item \emph{Say it:} \emph{māṇḍṇe}, once more
  \item \emph{Say it:} \emph{mīrāne baiṭhkīt āplaṁ mat māṇḍlaṁ.} — Meera put forward her view at the meeting.
  \item \emph{Recall:} say the Marathi for logic, then the Marathi for a reasoned argument, then say \emph{māṇḍṇe} again
\end{itemize}

\begin{tcolorbox}[breakable,colback=teal!4,colframe=teal!35!black,title={Before you move on}]
What does \mr{मांडणे} mean? (\textquotedblleft{}To put forward; to set out\textquotedblright{}.) Say it once more.
\end{tcolorbox}
\section[\texorpdfstring{\mr{पुरावा} (purāvā)}{पुरावा (purāvā)}]{\mr{पुरावा} (purāvā) — proof; evidence}

\label{lesson:MR-C343-purava}



\begin{quote}
Before the new one: say the Marathi for a reasoned argument, then the Marathi for to put forward.
\end{quote}

\subsection*{You'll want to know: \mr{पुरावा}}
\textbf{\mr{पुरावा}} — \emph{purāvā} — \textquotedblleft{}proof; evidence\textquotedblright{}.

\begin{quote}
\mr{तुमच्याकडे} \mr{याचा} \mr{काही} \mr{पुरावा} \mr{आहे} \mr{का}?
\end{quote}

\begin{quote}
\emph{tumacyākaḍe yācā kāhī purāvā āhe kā?}
\end{quote}

\begin{quote}
Do you have any proof of this?
\end{quote}

Masculine. Plural: \mr{पुरावे}, \emph{purāve}.

X \mr{चा} \mr{पुरावा}, proof of X; \mr{पुरावा} \mr{देणे}, to give evidence.

\begin{cousinweb}
\textbf{\mr{पुरावा}} (\emph{purāvā}) is a Marathi noun whose origin is not settled.
\end{cousinweb}

\begin{grammarlens}[title={What you've built}]
One more word for giving reasons and opinions.
\end{grammarlens}

\subsection*{Your turn}
\begin{itemize}
\raggedright
  \item \emph{Say it:} \emph{purāvā}
  \item \emph{Say it:} \emph{purāvā}, once more
  \item \emph{Say it:} \emph{tumacyākaḍe yācā kāhī purāvā āhe kā?} — Do you have any proof of this?
  \item \emph{Recall:} say the Marathi for a reasoned argument, then the Marathi for to put forward, then say \emph{purāvā} again
\end{itemize}

\begin{tcolorbox}[breakable,colback=teal!4,colframe=teal!35!black,title={Before you move on}]
What does \mr{पुरावा} mean? (\textquotedblleft{}Proof; evidence\textquotedblright{}.) Say it once more.
\end{tcolorbox}
